\section{Discussion}
\label{sec:discussion}

\noindent Our results should not be read as evidence that attention is irrelevant to grounding. IC, uncertainty baselines, and preliminary per-head probes all suggest that LVLMs contain object-relevant information, and some of that information may be attention-derived. The negative result is narrower and more actionable: raw attention mass, mean-over-head attention shape, and unselective attention amplification are not reliable certificates of grounding.

\noindent The practical implication is that attention-based claims need selectivity controls. A detector should show that its score separates present from absent objects at comparable generation positions and object categories. A mitigation method should show that its routing change improves true positives more than false positives, rather than merely increasing yes rate or caption object richness. Semantic-neighbor negatives make this requirement concrete by testing the cases where associated evidence is present but target evidence is absent.

\noindent The OWLv2 positive control makes the practical direction sharper. Detector-assisted grounding and counterfactual region tests already improve the semantic-neighbor slice because they ask whether the target category itself is supported, not merely whether the image is influential. The remaining method problem is to make that verification cheaper, better calibrated, and more integrated with decoding. Head selection and learned verification heads may still be useful, but they should be trained or selected against the same position, category, answer-prior, and semantic-neighbor controls before being interpreted as hallucination detection or mitigation.

\noindent A practical method should therefore satisfy four design requirements. First, it should score target evidence separately from visually associated evidence, since related-present negatives are where the proxy failure is most visible. Second, it should make selective decisions: a POPE intervention should reduce false positives without trading away comparable true positives, and a caption intervention should reduce hallucinated mentions without simply shortening captions or deleting object detail. Third, it should amortize expensive verification. The appendix TDEV-lite audit suggests one conservative path: supervised internal readouts cannot be used by themselves, but they can route a smaller set of high-risk cases into a stronger target verifier. Fourth, the verifier should eventually act at the point where claims are formed, for example through constrained regeneration or sentence-local correction, rather than only as a post-hoc score. These requirements define the method direction more than any particular OWLv2 backend: the contribution is the target-vs-neighbor verification criterion and its controls, not the external detector itself.
